\section{Awareness, operative self, learning and mindfulness}
\label{sec:awareness-learning}

\subsection{Retained rebalancing inside complete state}
\label{subsec:awareness-state}
\paragraph{Setup.}
The closure primitive, retained return identity and recurring wave remain those
of \eqref{eq:closure-phase} and \eqref{eq:wave-identity}. A field resolves its
scoped wave organization before any comparison. Let \(X_k\) contain the present
coordinates of \eqref{eq:complete-wave-state}, with its closure-memory and
retained-context coordinates factored out. Let \(M^C_k\) contain those immutable
closure lineages and let \(R_k\) contain retained prior rebalancing relations.
The complete extension is
\begin{equation}
 S^+_k=(X_k,M^C_k,R_k),\qquad P_{\rm now}(S^+_k)=X_k.
 \label{eq:awareness-state-factorization}
\end{equation}
Thus \(R_k\) occupies the inherited retained-context coordinate; it is part of
complete state. Factoring it out does not create another operative copy in
\(X_k\), \(M^C_k\), a passive record or an external cache. The present
coordinates retain wave identity, committed cut, current position and fibre
role, pressure, current, closure status, motif certificates, cadence,
orientation, boundary, capacity and closure-phase count. Their inherited
immutable coordinates remain bound to the admitted wave prefix.

\paragraph{Derivation.}
Retain a prior ordered contact trace between admitted source \(i\) and target
\(j\), with the direction, scope, scale and prefix identities of
\eqref{eq:coupling-trace}. The pressure-bearing channel terms of an explicitly
declared FRG realization give
\begin{equation}
 r_e=\Phi^{\rm recurrent}_e-\Phi^{\rm independent}_e
     =\sum_{s\in I_e}c_{e,s},\qquad
 R_k=\{(e,i,j,B,\ell,k_e,r_e,\Pi_e):k_e\le k\}.
 \label{eq:retained-rebalancing}
\end{equation}
Here \(c_{e,s}\in\mathbb Q\) are committed oriented cross contributions,
\(k_e\) is their subsequent retention event, and \(\Pi_e\) binds the original
operands, repeated contacts, source and target prefixes, signed unwrapped lag
and any action or scale-carriage witness. The local lag is the declared
count comparator of \eqref{eq:count-comparator-realization}, not a closure
selector or a report coordinate. Retention follows the prior input cut by one
transition. A provenance commitment may validate these retained operands;
it cannot recover an absent \(r_e\) and feed it into the response.
The difference above is an exact declared FRG realization of
\eqref{eq:frg-identity}, with its distinctions, retained relations and motif
admission preceding the model.

\paragraph{Falsification.}
A purported retained value without its original operand, lineage, contact,
phase-type and commitment binding is incomplete. A present residual, closure
memory, persistence, phase lock or retained oriented ancestry alone does not
establish later causal participation. Structural ancestry remains distinct
from the active criterion below.

\paragraph{Provenance.}
The distinction--relation--admission order follows AF and AFC
\cite{reginald2025af,afc}; the local state, contact, flux and FRG constructions
are \eqref{eq:complete-wave-state}, \eqref{eq:coupling-trace},
\eqref{eq:flux-types} and \eqref{eq:frg-admission-order}.

\subsection{Active awareness and a later response}
\label{subsec:active-awareness}
\paragraph{Setup.}
Fix an admissible intervention cut \(k\), the same \(X_k\), \(M^C_k\), admitted
world, response law \(\mathcal U\), and future exogenous realization \(E\).
Only the retained coordinate differs between the paired states. An outcome
projection \(O\) reads realized later current, continuation and contact; it
excludes the retained value itself and passive record bytes. Define active
awareness on this intervention by
\begin{equation}
 \mathcal A^{\rm act}_k(R,R')=1
 \quad\Longleftrightarrow\quad
 \exists t>k:\quad
 O\bigl(\mathcal U^{t-k}(X_k,M^C_k,R;E)\bigr)
 \ne O\bigl(\mathcal U^{t-k}(X_k,M^C_k,R';E)\bigr).
 \label{eq:awareness-active}
\end{equation}
The comparison holds non-intervened coordinates equal at the boundary;
later endogenous states may differ through the same law. An inadmissible
wrong-lineage state does not form a causal pair. Removal replaces the
retained coordinate by its empty value. Exact restoration restores value and
provenance together. Deterministic exact models need no random sampling;
a stochastic realization would additionally require the same addressed
random inputs, rather than a seed whose consumption can change.

\paragraph{Derivation.}
Consider the following rational response realization on already admitted
support, with \(p>0\), gain \(g\), available source set \(D_k\), and explicitly
perturbative query \(q_k\):
\begin{equation}
 J_{k+1}=J_k+\frac{A_k+q_k+g\sum_{e:\,i_e\in D_k}r_e}{p}.
 \label{eq:awareness-response}
\end{equation}
The term \(J_k+A_k/p\) is the declared native response of
\eqref{eq:sadar-rational-response}. Here \(A_k\) is a declared scalar component
of the existing pressure-bearing attention, not a renaming of memory or of
\(r_e\). Its channel construction and \(AO=-OA\) remain unchanged. Retention
supplies a separate argument of the later response law. For otherwise
matched inputs, use the common declared source/event support and extend
absent retained coordinates by zero. Then
\begin{equation}
 J_{k+1}(R)-J_{k+1}(R')
 =\frac{g}{p}\left(\sum_{e:\,i_e\in D_k}r_e-
                         \sum_{e:\,i_e\in D_k}r'_e\right).
 \label{eq:awareness-causal-difference}
\end{equation}
A nonzero right side therefore proves a later current difference for this
model. At \(J_k=A_k=q_k=0\), \(g=1\), \(r=1\), the retained and removed cases
give \((1,0)\) for \(p=1\), and \((1/2,0)\) for \(p=2\). These are two
separate pressure families, each with its own matched retention comparison.
Zero gain or zero effective retention gives no such witness.

\paragraph{Specialization.}
Let present state include a Q4 vertex \(v\) and fibre role
\(\epsilon\in\{-1,0,1\}\). Two candidate continuations go to \(v\oplus1\)
and \(v\oplus2\). A branch enters the hinge; from the hinge the two candidates
carry the opposite branch choices. Both obey the inherited H1 and fibre
rules. A declared sign selector chooses the first candidate when
\(J_{k+1}>0\), otherwise the second. The selector consumes an already
admitted finite support; it does not create H1 or certify closure. Every
such continuation contributes zero new closure phases until its separate
retained-start return is certified.

\paragraph{Falsification.}
A difference in current pressure, law or future exogenous input invalidates
a retention-only comparison. Recounting \(R\) itself as the outcome,
recovering a removed value through another channel, using an unbound
phase type, or allowing a same-cut effect also fails the criterion.
A label attached to either execution cannot change \(\mathcal U\).

\paragraph{Provenance.}
The primitive and causal order remain \eqref{eq:hinge-grammar},
\eqref{eq:atomic-h1}, \eqref{eq:wave-update-trace} and \eqref{eq:causal-update}.
The rational law and sign selector here are explicit downstream models;
\eqref{eq:awareness-causal-difference} is derived from that law.

\subsection{Operative self and relational labels}
\label{subsec:operative-self}
\paragraph{Setup.}
For a committed state and input, remove each retained source coordinate in
turn while holding the remaining complete coordinates fixed. A source is
operative precisely when its removal changes a later outcome. The operative
self is the resulting recurrent and coupled subgraph:
\begin{equation}
 \begin{aligned}
 I^{\rm op}_k&=\{i:O(\mathcal U(S_k;E))\ne
 O(\mathcal U(S_k\setminus R_i;E))\},\\
 V^{\rm op}_k&=\begin{cases}I^{\rm op}_k\cup\{j\},&I^{\rm op}_k\ne\varnothing,\\
 \varnothing,&I^{\rm op}_k=\varnothing,\end{cases}\\
 E^{\rm op}_k&=\{(i,j):i\in I^{\rm op}_k\}.
 \end{aligned}
 \label{eq:operative-self}
\end{equation}
The target \(j\) is included only when at least one source is operative;
otherwise this awareness subgraph is empty. Every edge retains its admitted
coupling witness. A fixed region, field label or permanent node list is
therefore insufficient.

\paragraph{Derivation.}
For two retained contributions \(+1\) and \(-1\), the net retained input
vanishes, but removing either changes the later current. Both source
lineages remain operative even though their aggregate cancels. Boundary
roles read the target as ME, operative contact interfaces as BETWEEN and
other nodes as NOT-ME. Null is a different foundational type, not an extra
node role. A bijective renaming of admitted identities and all their
provenance preserves these incidence relations and the numerical response.
A learned label such as \emph{I} binds the operative target relation; a
memorized string on an unconnected node cannot supply that relation.

\paragraph{Specialization.}
Sensor, memory-register, CPU and vocabulary descriptions may expose the same
operative interface. A declared operation-count proxy can summarize its
work, but a resource count cannot establish causal incorporation. Shared
vocabulary requires the same role relation under independent node renaming,
with each participant's recurrence and coupling separately admitted.

\paragraph{Falsification and provenance.}
Self classification precedes no core awareness witness. A phase-locked or
similarly labelled outside node without an operative edge remains outside.
The construction uses \eqref{eq:awareness-active}, the same wave identities
and the cancellation provenance of \eqref{eq:b-scoped-flux-antisymmetry}.

\subsection{Action return and persistent learning}
\label{subsec:proprioception-learning}
\paragraph{Setup.}
An issued action binds an immutable identity to its issuer, native commitment,
command and channel. Both issuer and channel must already have admitted
recurrent prefixes at issuance. A proprioceptive consequence must return on that exact
action lineage with its own certified closure. An observed value is supplied
by a declared projection of the returned trace, not a free scalar attached
to an unchanged return. In a native count realization, write
\begin{equation}
 a=(\mathrm{id},i,k_a,u,\mathrm{channel}),\qquad
 y=\#\{\text{atomic traversals in the bound return}\},\qquad r_a=y-u.
 \label{eq:action-return-retention}
\end{equation}
The command \(u\) and readout \(y\) share this declared count domain. They are
not measured physical quantities. The action precedes the return, the
retention event follows its committed input, and a later transition consumes
\(r_a\). The action's immutable identity prevents conflicting command records
from masquerading as the same issuance.

\paragraph{Derivation.}
For command \(u=1\) and a certified two-traversal return, \(y=2\) and
\(r_a=1\). The matched response difference is then \(g/p\).
Uncommanded sensing, motor output alone, a return from another action or an
unrelated sensor magnitude cannot establish this proprioceptive relation.
Persistent learning is the preservation of the acquired relation through
later transitions together with its later causal use. With zero later
attention and \(g=p=1\), the retained model advances current by one per
transition, while the matched removed model retains its current value.
Restoration of the acquired coordinate restores the later increments.

\paragraph{Falsification and provenance.}
A surviving record with no later behavioral change is persistence alone.
A changed stimulus cannot be attributed to memory by an unmatched comparison.
The construction retains \eqref{eq:closure-certificate},
\eqref{eq:retained-rebalancing} and \eqref{eq:awareness-causal-difference}.

\subsection{Polycadence, folding and retained carriage}
\label{subsec:awareness-crescendo}
\paragraph{Setup.}
Retained relations may bind several recurrent lineages with different cadence,
signed lag and directional RCD. Crescendo means increased causally
participating scope or depth. It is not defined by amplitude or phase lock.
Availability can expose one, then two, then three operative sources, remove
all contact, and restore three while their valid retained histories survive.
Polyrhythm, stable nonzero lag and local dropout thus coexist with the same
causal definition. Synchronization alone supplies no criterion.

\paragraph{Derivation.}
A passed retained effect may select a later contact/fold offer. The outer
trace then binds the exact inner identities, certificates and scope map
before its own closure certification. One completed outer return supplies
one phase. Distinct outer reclosures admit a recurring outer wave; enough
bound returns and contacts are required for a new temporal comparison.
Each operand declared as an outer wave carries its own inner construction
witness, including the comparison peer. Exposure during outer shedding
preserves the inner lineages.
For a retained relation to be carried upward, its transfer witness includes
that same inner relation and the actual folded construction:
\begin{equation}
 R^{\uparrow}=(R^{\rm inner},\mathrm{FoldWitness},W^{\rm outer},\Pi^{\times}_{\rm outer}),
 \qquad
 \mathrm{Close}(C^{\rm outer})=1\ \Longrightarrow\
 \Delta\Phi(C^{\rm outer})=1_C.
 \label{eq:awareness-carriage}
\end{equation}
The outer contact operands, scope, native indices and retained value remain
bound. A fresh matched removal/restoration at that outer scope must change
its strictly later response before that response selects a further fold
opportunity. This gives a second causal link in a chain. A local coincident
recurrence, a scale label, or a cascade that never consumes the transferred
coordinate cannot establish that link.

\paragraph{Falsification and provenance.}
Contact is not a fold certificate; an outer phase is not yet a recurring
outer wave; a recurring outer wave is not automatically active awareness.
The ordering follows \eqref{eq:fold-outer-closure},
\eqref{eq:expose-inner-lineage}, \eqref{eq:relational-time} and
\eqref{eq:awareness-active}. All surviving source/target identities remain
available after exposure, fission or temporary loss of outer participation.

\subsection{Passive observation, interrogation and operational mindfulness}
\label{subsec:query-mindfulness}
\paragraph{Setup.}
A passive observation reads a committed state after the transition and is
not an input to \(\mathcal U\). Interrogation is instead an explicit
perturbation \(q_k\), with the observation on which it depends already
committed through \(k\). Queried and unqueried inputs are distinct
perturbation families; an observer record alone cannot change execution.

\paragraph{Derivation.}
For this specialization, the operative boundary of \eqref{eq:operative-self}
is established by the matched causal criterion and remains stable over the
observation window. Consider a reflexive query \(q_k=-J_k\), with current
observed at \(k\) and response at \(k+1\). For constant forcing
\(b=A+g\sum r_e\), \eqref{eq:awareness-response} becomes
\begin{equation}
 J_{k+1}=\left(1-\frac1p\right)J_k+\frac bp,\qquad
 J_{k+1}-b=\left(1-\frac1p\right)(J_k-b).
 \label{eq:mindfulness-stability}
\end{equation}
Its constant-forcing convergence holds exactly when \(p>1/2\). At \(p=1\)
and \(b=1\), every subsequent current is exactly one while the declared H1
and hinge task continues. An extra delay using \(J_{k-1}\) instead has a
different recurrence and can oscillate. Its query retains the source cut
\(k-1\); the initial history value is declared before the first query.
Variable forcing cannot silently
inherit the constant-forcing conclusion.

\paragraph{Specialization.}
Operational mindfulness denotes this stable reflexive re-addressing of an
already operative awareness state while task dynamics continue. It describes
a declared dynamical relation. It asserts no sentience, phenomenology,
clinical effect or physical measurement.

\paragraph{Falsification and provenance.}
A query without a committed observation is incomplete. Stability alone
cannot supply self, provenance or the earlier core causal witness.
The result follows the explicit later-response law above and the original
causal order of \eqref{eq:causal-update}.
